%----------------------------------------------------------------------------------------
% DE AFGELEIDE. Alle opgaven zijn eigen werk.
%
% Opbouw: Start, leerstof (theorie, werkwijze of uitleg, voorbeeld, korte oefening), en op een
% nieuwe pagina de Opdrachten (door elkaar, stijgende moeilijkheid), Koppel en Test jezelf.
% Labels van dit hoofdstuk beginnen met h2-.
%----------------------------------------------------------------------------------------

\mychapter{2}{De afgeleide}

\lkr{Leerplan B+S: LPD 18-21 (afgeleiden), het deel over het differentiequotiënt, de afgeleide als limiet, de raaklijn, de afgeleide functie en de rekenregels. Het beeld van de voorkant (koorden door $P$ die naar de raaklijn draaien) is de figuur in sectie 1. Laat leerlingen de definitie een paar keer zelf uitwerken voor ze de rekenregels gebruiken.}

Een snelheidsmeter toont hoe snel je \textit{nu} rijdt, niet gemiddeld over de rit. Maar snelheid is afstand gedeeld door tijd, en op één ogenblik leg je geen afstand af. Hoe kan dat dan? Het antwoord is een limiet: neem de gemiddelde snelheid over een steeds kortere tijd. Zo ontstaat de afgeleide, misschien het belangrijkste begrip van de analyse.

\begin{oefeningen}[Start]
\oef{h2-start-1}

\oef{h2-verken}
\end{oefeningen}

%========================================================================================
\section{De afgeleide in een punt}

\naast{%
De rechte door twee punten $P(a, f(a))$ en $Q(a + h, f(a + h))$ van de grafiek is een \textbf{koorde}. Haar richtingscoëfficiënt is het differentiequotiënt: de gemiddelde verandering van $f$ tussen $a$ en $a + h$.

Laat $Q$ naar $P$ schuiven ($h \to 0$). De koorden draaien dan naar een grensstand: de \textbf{raaklijn} in $P$.}{%
\begin{tikzpicture}[x=1.1cm,y=.55cm,font=\scriptsize]
  \draw[->,lijngrijs] (-.3,0) -- (2.6,0) node[right,text=black] {$x$};
  \draw[->,lijngrijs] (0,-.3) -- (0,5.6) node[above,text=black] {$y$};
  \kromme{0}{2.3}{(\x)^2}
  \draw[grijs] (.55,-.35) -- (2.3,4.9);
  \draw[grijs,dashed] (.5,-.25) -- (2.3,4.25);
  \draw[accent,thick] (.3,-.4) -- (2.2,3.4);
  \fill[accent] (1,1) circle (.05) node[above left,text=black] {$P$};
  \fill (2,4) circle (.05) node[right] {$Q$};
  \fill (1.5,2.25) circle (.04);
  \node[accent,right] at (2.2,3.3) {raaklijn};
\end{tikzpicture}}

\begin{theorie}
Het \textbf{\textcolor{accent}{differentiequotiënt}} van $f$ tussen $a$ en $a + h$ is $\formule{\frac{f(a + h) - f(a)}{h}}$. De \textbf{\textcolor{accent}{afgeleide}} van $f$ in $a$ is
\[\formule{f'(a) = \lim_{h \to 0} \frac{f(a + h) - f(a)}{h}},\]
als die limiet bestaat; $f$ is dan \textbf{\textcolor{accent}{afleidbaar}} in $a$. Het getal $f'(a)$ is de richtingscoëfficiënt van de raaklijn in $(a, f(a))$:
\[\formule{y = f(a) + f'(a)(x - a)}.\]
\end{theorie}

\voorbeeld{}{Bereken met de definitie $f'(3)$ voor $f(x) = x^2$, en geef de raaklijn in $(3, 9)$.\label{vb:h2-definitie}}{%
\uitleg{Werk het differentiequotiënt uit tot de $h$ in de noemer wegvalt. Neem dan de limiet.}
\begin{stappenbreed}
\frac{(3 + h)^2 - 9}{h} = \frac{6h + h^2}{h} = 6 + h & Voor $h \neq 0$. \\
f'(3) = \lim_{h \to 0} (6 + h) = 6 & \\
y = 9 + 6(x - 3) = 6x - 9 & Raaklijn.
\end{stappenbreed}}

\begin{oefeningen}
\oef{h2-defpunt}
\end{oefeningen}

Is $s(t)$ de afgelegde weg op tijdstip $t$, dan is het differentiequotiënt de gemiddelde snelheid, en $s'(t)$ de \textbf{ogenblikkelijke snelheid}.

\voorbeeld{}{Een steen valt; na $t$ seconden heeft hij $s(t) = 5t^2$ meter afgelegd. Hoe snel valt hij na 2 seconden?\label{vb:h2-snelheid}}{%
\uitleg{Neem de gemiddelde snelheid tussen $2$ en $2 + h$, en laat $h$ naar 0 gaan.}
\begin{stappenbreed}
\frac{5(2 + h)^2 - 20}{h} = \frac{20h + 5h^2}{h} = 20 + 5h & Gemiddelde snelheid. \\
s'(2) = \lim_{h \to 0} (20 + 5h) = 20 & In m/s, ongeveer 72 km/h.
\end{stappenbreed}}

\begin{oefeningen}
\oef{h2-definitie}
\end{oefeningen}

%========================================================================================
\section{De afgeleide functie}

Bereken je de afgeleide in elk punt $x$, dan krijg je een nieuwe functie.

\begin{theorie}
De \textbf{\textcolor{accent}{afgeleide functie}} van $f$ is $\formule{f'(x) = \lim_{h \to 0} \frac{f(x + h) - f(x)}{h}}$. Andere notaties: $\frac{\mathrm{d}f}{\mathrm{d}x}$ of $\frac{\mathrm{d}y}{\mathrm{d}x}$.
\begin{stelling}{Afgeleide van een macht}
Voor elk rationaal getal $r$: \quad $\formule{(x^r)' = r\,x^{r - 1}}$ \quad (waar $x^r$ en $x^{r-1}$ bestaan).
\end{stelling}
\begin{bewijs}
Voor een natuurlijk getal $n$: met het binomium van Newton is
\bewijsregel{(x + h)^n = x^n + n x^{n-1} h + h^2 \cdot (\ldots),}
dus $\frac{(x + h)^n - x^n}{h} = n x^{n-1} + h \cdot (\ldots) \to n x^{n-1}$. Voor andere exponenten: zie de voorbeelden en opdrachten.
\end{bewijs}
\end{theorie}

Zo is $(x^3)' = 3x^2$, $\bigl(\frac1x\bigr)' = (x^{-1})' = -x^{-2} = -\frac{1}{x^2}$ en $(\sqrt x)' = \bigl(x^{\frac12}\bigr)' = \frac12 x^{-\frac12} = \frac{1}{2\sqrt x}$. De afgeleide van een constante is $0$.

\voorbeeld{}{Bewijs met de definitie dat $\bigl(\frac1x\bigr)' = -\frac{1}{x^2}$.\label{vb:h2-functie}}{%
\uitleg{Zet de teller van het differentiequotiënt op één noemer.}
\begin{stappenbreed}
\frac{\frac{1}{x + h} - \frac1x}{h} = \frac{x - (x + h)}{h\,x(x + h)} = \frac{-1}{x(x + h)} & Op één noemer, $h$ schrappen. \\
\lim_{h \to 0} \frac{-1}{x(x + h)} = -\frac{1}{x^2} &
\end{stappenbreed}}

\begin{oefeningen}
\oef{h2-deffunctie}

\oef{h2-functie}
\end{oefeningen}

%========================================================================================
\section{Rekenregels}

\begin{theorie}
\begin{stelling}{Rekenregels voor afgeleiden}
Voor afleidbare functies $f$ en $g$ en een getal $c$:
\[\formule{(cf)' = c f'} \qquad \formule{(f + g)' = f' + g'} \qquad \formule{(fg)' = f'g + fg'} \qquad \formule{\Bigl(\frac{f}{g}\Bigr)' = \frac{f'g - fg'}{g^2}}\]
\end{stelling}
\begin{bewijs}
\textit{Productregel.} Tel $f(x)g(x + h)$ op en trek het af:
\bewijsregel{f(x + h)g(x + h) - f(x)g(x) = \bigl(f(x + h) - f(x)\bigr)g(x + h) + f(x)\bigl(g(x + h) - g(x)\bigr).}
Deel door $h$ en laat $h \to 0$: de eerste term gaat naar $f'(x)g(x)$ (want $g$ is continu), de tweede naar $f(x)g'(x)$.
\end{bewijs}
\end{theorie}

\voorbeeld{}{Bereken de afgeleide van $f(x) = 3x^4 - 2\sqrt x + \dfrac5x$.\label{vb:h2-som}}{%
\uitleg{Schrijf elke term als macht en gebruik de som- en veelvoudregel.}
\begin{stappenbreed}
f(x) = 3x^4 - 2x^{\frac12} + 5x^{-1} & Als machten. \\
f'(x) = 12x^3 - x^{-\frac12} - 5x^{-2} = 12x^3 - \frac{1}{\sqrt x} - \frac{5}{x^2} &
\end{stappenbreed}}

\begin{oefeningen}
\oef{h2-som}
\end{oefeningen}

\voorbeeld{}{Bereken de afgeleide van $f(x) = (x^2 + 1)(x^3 - x)$ met de productregel. Controleer door eerst uit te werken.\label{vb:h2-product}}{%
\uitleg{Eerste factor afgeleid maal de tweede, plus de eerste maal de tweede afgeleid.}
\begin{stappenbreed}
f'(x) = 2x(x^3 - x) + (x^2 + 1)(3x^2 - 1) & Productregel. \\
= 2x^4 - 2x^2 + 3x^4 + 2x^2 - 1 = 5x^4 - 1 & \\
f(x) = x^5 - x \Rightarrow f'(x) = 5x^4 - 1 & Controle.
\end{stappenbreed}}

\begin{oefeningen}
\oef{h2-product}
\end{oefeningen}

\voorbeeld{}{Bereken de afgeleide van $f(x) = \dfrac{2x + 1}{x - 3}$.\label{vb:h2-quotient}}{%
\uitleg{Quotiëntregel: (teller afgeleid maal noemer min teller maal noemer afgeleid) gedeeld door de noemer in het kwadraat.}
\begin{stappenbreed}
f'(x) = \frac{2(x - 3) - (2x + 1) \cdot 1}{(x - 3)^2} = \frac{-7}{(x - 3)^2} & Quotiëntregel.
\end{stappenbreed}
De afgeleide is overal negatief: de functie daalt op elk interval van haar domein.}

\begin{oefeningen}
\oef{h2-quotient}

\oef{h2-rekenregels}
\end{oefeningen}

%========================================================================================
\section{Afleidbaar en continu}

Een functie is niet overal afleidbaar. Bij $f(x) = |x|$ is het differentiequotiënt in $0$ gelijk aan $\frac{|h|}{h}$: $-1$ voor $h < 0$ en $1$ voor $h > 0$. De limiet bestaat niet: de grafiek heeft een knik. Omgekeerd geldt wel: een functie die afleidbaar is in $a$, is daar ook continu. Een sprong in de grafiek maakt de afgeleide dus zeker onmogelijk.

\begin{oefeningen}
\oef{h2-knik}
\end{oefeningen}

%========================================================================================
\opdrachtenkop

\begin{oefeningen}
\needspace{13\baselineskip}
\oef{h2-o10}

\needspace{9\baselineskip}
\oef{h2-o13}

\needspace{8\baselineskip}
\oef{h2-o2}

\needspace{9\baselineskip}
\oef{h2-o16}

\needspace{10\baselineskip}
\oef{h2-o1}

\needspace{10\baselineskip}
\oef{h2-o14}

\oef{h2-o12}

\oef{h2-o15}

\needspace{8\baselineskip}
\oef{h2-o3}

\needspace{12\baselineskip}
\oef{h2-o4}

\needspace{9\baselineskip}
\oef{h2-o23}

\needspace{14\baselineskip}
\oef{h2-o18}

\oef{h2-o5}

\oef{h2-o20}

\needspace{15\baselineskip}
\oef{h2-o11}

\needspace{12\baselineskip}
\oef{h2-o17}

\oef{h2-o6}

\needspace{12\baselineskip}
\oef{h2-o24}

\needspace{9\baselineskip}
\oef{h2-o19}

\oef{h2-o7}

\oef{h2-o21}

\needspace{9\baselineskip}
\oef{h2-o27}

\oef{h2-o8}

\needspace{10\baselineskip}
\oef{h2-o28}

\needspace{10\baselineskip}
\oef{h2-o26}

\oef{h2-o9}

\oef{h2-o29}
\end{oefeningen}

%----------------------------------------------------------------------------------------
% KOPPEL: verbanden met andere leerstof
%----------------------------------------------------------------------------------------
\needspace{20\baselineskip}
\opdrachtsectie{Koppel}

\lkr{Integralen en Rationale functies: kijk na wat al gezien is. Het pakket legt de volgorde niet vast.}

\begin{oefeningen}
\oef{h2-koppel-1}

\oef{h2-koppel-2}

\end{oefeningen}

%----------------------------------------------------------------------------------------
\opdrachtsectie{Test jezelf}

\begin{oefeningen}
\oef{h2-t1}

\oef{h2-t2}

\oef{h2-t3}

\oef{h2-t4}

\oef{h2-t5}

\oef{h2-t6}

\oef{h2-t7}

\oef{h2-t8}
\end{oefeningen}

\herhaalschema{%
  \fv{oef:h2-t1}{vb:h2-definitie}{oef:h2-defpunt}%
  \fv{oef:h2-t2}{vb:h2-definitie}{oef:h2-o2}%
  \fv{oef:h2-t3}{vb:h2-som}{oef:h2-som}%
  \fv{oef:h2-t4}{vb:h2-product}{oef:h2-product}%
  \fv{oef:h2-t5}{vb:h2-quotient}{oef:h2-quotient}%
  \fv{oef:h2-t6}{vb:h2-snelheid}{oef:h2-o15}%
  \fv{oef:h2-t7}{vb:h2-functie}{oef:h2-deffunctie}%
  \fzv{oef:h2-t8}{oef:h2-knik}}

\ifzelfstudie
\section*{Maak je eigen samenvatting}
Sluit dit boek en schrijf in je schrift: het differentiequotiënt en zijn betekenis (koorde, gemiddelde snelheid), de definitie van $f'(a)$ en van de afgeleide functie, de vergelijking van de raaklijn, de afgeleide van $x^r$, de rekenregels met het bewijs van de productregel, en een voorbeeld van een functie die niet afleidbaar is. Vergelijk daarna met de kaders.
\fi
